% ---------------------------------------------------------------------------------------------------------------
% Re-evaluation of the Chapter 4 results from the saved checkpoints and training logs (results/08_thesis_recheck).
% Can be used as an appendix section, or its paragraphs can replace the affected text of Sections 4.2.1 / 4.2.2.
% Requires: graphicx, booktabs, adjustbox. \ResultsDir must point to the results/ directory.
% ---------------------------------------------------------------------------------------------------------------

\section{Re-evaluation of the Chapter~4 Results}
\label{sec:recheck}

The saved checkpoints of Chapter~4 were evaluated again with the thesis' own evaluation pipeline: the ClimateNet test
images with prompts derived from the ground truth (box, point or a random choice per image), \texttt{ClimateSAM.\allowbreak infer},
and TC and AR scored as two independent binary problems. Available were the four Infused Token models of the appendix
table (linear and nonlinear input adapter, MLP ratio 1.0 and 0.5) and all Phase-2 prompter checkpoints; the Concat
Token, Single Token and LoRA checkpoints are not available, and for Single Token and Single LoRA the training logs were
used instead.

\subsection{The Phase-1 validation skips the input adapter}

The validation routine of the Phase-1 training script encodes the images with \texttt{ClimateSAM.\allowbreak set\_\allowbreak infer\_\allowbreak img()},
which passes the first three raw channels (TMQ, U850, V850) to the encoder and does not apply the learned input adapter;
training uses \texttt{encode\_\allowbreak images()}, which does. Table~\ref{tab:thesis_recheck_phase1} and
Figure~\ref{fig:recheck_phase1} compare the reported values with a re-run of the original evaluation and with a re-run in
which the adapter is applied.

\begin{table}[htbp]
\centering
\small
\begin{adjustbox}{max width=\linewidth}
\begin{tabular}{llccccccc}
\toprule
Infused Token & Prompt & Thesis TC & Thesis AR & As in thesis TC & As in thesis AR & Adapter applied TC & Adapter applied AR & $\Delta$ FG (applied $-$ thesis) \\
\midrule
linear, MLP 1.0 & box & 0.724 & 0.632 & 0.730 & 0.616 & 0.736 & 0.625 & +0.002 \\
 & point & 0.522 & 0.456 & 0.525 & 0.442 & 0.532 & 0.453 & +0.003 \\
 & random & 0.595 & 0.523 & 0.624 (0.613--0.642) & 0.521 (0.521--0.521) & 0.637 (0.622--0.657) & 0.530 (0.529--0.532) & +0.025 \\
\midrule
linear, MLP 0.5 & box & 0.724 & 0.645 & 0.738 & 0.648 & 0.741 & 0.650 & +0.011 \\
 & point & 0.526 & 0.469 & 0.526 & 0.470 & 0.533 & 0.486 & +0.012 \\
 & random & 0.606 & 0.536 & 0.627 (0.612--0.645) & 0.555 (0.554--0.556) & 0.636 (0.621--0.653) & 0.560 (0.558--0.561) & +0.027 \\
\midrule
nonlinear, MLP 1.0 & box & 0.699 & 0.445 & 0.733 & 0.439 & 0.721 & 0.556 & +0.067 \\
 & point & 0.331 & 0.313 & 0.394 & 0.314 & 0.520 & 0.440 & +0.158 \\
 & random & 0.447 & 0.360 & 0.523 (0.513--0.541) & 0.356 (0.349--0.364) & 0.623 (0.606--0.643) & 0.491 (0.486--0.496) & +0.153 \\
\midrule
nonlinear, MLP 0.5 & box & 0.733 & 0.522 & 0.749 & 0.536 & 0.748 & 0.622 & +0.057 \\
 & point & 0.476 & 0.351 & 0.461 & 0.373 & 0.531 & 0.462 & +0.083 \\
 & random & 0.572 & 0.421 & 0.613 (0.591--0.638) & 0.442 (0.439--0.447) & 0.642 (0.617--0.673) & 0.530 (0.528--0.534) & +0.090 \\
\bottomrule
\end{tabular}
\end{adjustbox}
\caption{Re-evaluation of the Phase-1 checkpoints (appendix Table 1) on the 61 test images with the thesis pipeline. "As in thesis": images encoded with set\_infer\_img (first three raw channels, learned input adapter skipped), as in the validation of train\_adaptation.py; "adapter applied": encode\_images, as in training. Random prompts: mean of three prompt draws (range).}
\label{tab:thesis_recheck_phase1}
\end{table}


\begin{figure}[htbp]
\centering
\includegraphics[width=\linewidth]{\ResultsDir/figures/thesis_recheck_phase1.pdf}
\caption{Phase-1 checkpoints with ground-truth prompts: value reported in Table~4.12 / appendix Table~1 (ring),
re-run with the original evaluation (grey, input adapter skipped) and with the input adapter applied (red). Random
prompts: mean and range over three prompt draws.}
\label{fig:recheck_phase1}
\end{figure}

For the linear adapter the reported box and point values reproduce within about $\pm0.015$ IoU, and applying the adapter
changes them by at most 0.02: the adapter keeps its identity weights for TMQ, U850 and V850 (Table~4.7), so the raw
channels have the same spatial structure as its output. The nonlinear adapter, in contrast, maps the inputs to
low-contrast images (Figure~\ref{fig:recheck_adapter_inputs}); its encoder was trained on these images but evaluated on
the raw channels. With the adapter applied, the nonlinear models gain 0.04--0.13 mean FG IoU and 0.09--0.14 AR IoU, and
the advantage of the linear adapter shrinks from 0.06--0.17 to 0.005--0.042 mean FG IoU
(Table~\ref{tab:thesis_recheck_adapter_bootstrap}), with the remaining gain almost entirely in AR. The large gap of
Table~4.6 is therefore caused by the evaluation and not by spatial blurring in the nonlinear adapter.
Figure~\ref{fig:recheck_maps_nonlinear} shows the effect on one test image.

\input{\ResultsDir/tables/thesis_recheck_adapter_bootstrap}

\begin{figure}[htbp]
\centering
\includegraphics[width=\linewidth]{\ResultsDir/figures/thesis_recheck_adapter_inputs.pdf}
\caption{Input of the image encoder for test image 17: the raw channels used at evaluation (a--c) and the outputs of
the linear (d--f) and nonlinear (g--i) input adapters on which the respective encoders were trained (grey scale 0--255).}
\label{fig:recheck_adapter_inputs}
\end{figure}

\begin{figure}[htbp]
\centering
\includegraphics[width=\linewidth]{\ResultsDir/figures/thesis_recheck_maps_nonlinear_17.pdf}
\caption{Infused Token (MLP ratio 1.0) with ground-truth box prompts on test image 17: linear adapter as evaluated in
the thesis (b), nonlinear adapter as evaluated in the thesis (c) and with its adapter applied (d).}
\label{fig:recheck_maps_nonlinear}
\end{figure}

\subsection{Checkpoint selection on the test set}

The Phase-1 training script validates on the 61 test images every five epochs and saves the epoch with the best mean
of TC and AR IoU. Figure~\ref{fig:recheck_training_logs} shows the logged test IoU of the training runs. For the
token-based adapters the curves plateau and the selected epoch differs from the last one by at most 0.01 IoU. For
Single LoRA the TC IoU collapses during training, from 0.63 at epoch~6 to 0.10 at epoch~50 for rank~32 and from 0.59 to
0.12 for rank~64; the LoRA values of Tables~4.6 and~4.8 therefore describe the best early test epoch, and both ranks,
not only rank~64, are unstable for TCs.

\begin{table}[htbp]
\centering
\small
\begin{adjustbox}{max width=\linewidth}
\begin{tabular}{lcccc}
\toprule
Run (wandb) & Max TC IoU & Max AR IoU & Best mean epoch TC / AR & Last epoch TC / AR \\
\midrule
infused\_token\_vit\_b\_mlp1\_CORRECTED & 0.648 (ep. 46) & 0.558 (ep. 46) & 0.648 / 0.558 (ep. 46) & 0.641 / 0.556 (ep. 50) \\
infused\_token\_vit\_b\_mlp05\_CORRECTED & 0.647 (ep. 46) & 0.595 (ep. 46) & 0.647 / 0.595 (ep. 46) & 0.644 / 0.594 (ep. 50) \\
single\_sam\_vit\_b\_CORRECTED & 0.644 (ep. 46) & 0.548 (ep. 50) & 0.644 / 0.546 (ep. 46) & 0.641 / 0.548 (ep. 50) \\
lora\_single\_vit\_b\_r32\_CORRECTED & 0.634 (ep. 6) & 0.594 (ep. 21) & 0.634 / 0.568 (ep. 6) & 0.104 / 0.579 (ep. 50) \\
lora\_single\_vit\_b\_r64\_CORRECTED & 0.592 (ep. 11) & 0.591 (ep. 16) & 0.592 / 0.569 (ep. 11) & 0.116 / 0.588 (ep. 50) \\
infused\_token\_vit\_b\_mlp1\_nonlinear & 0.478 (ep. 46) & 0.352 (ep. 41) & 0.477 / 0.352 (ep. 41) & 0.477 / 0.344 (ep. 50) \\
infused\_token\_vit\_b\_mlp05\_nonlinear & 0.586 (ep. 50) & 0.424 (ep. 46) & 0.586 / 0.424 (ep. 46) & 0.586 / 0.422 (ep. 50) \\
single\_sam\_vit\_b\_nonlinear & 0.589 (ep. 46) & 0.451 (ep. 36) & 0.587 / 0.449 (ep. 41) & 0.589 / 0.447 (ep. 50) \\
lora\_single\_vit\_b\_r32\_nonlinear & 0.481 (ep. 11) & 0.506 (ep. 16) & 0.481 / 0.458 (ep. 11) & 0.034 / 0.496 (ep. 50) \\
lora\_single\_vit\_b\_r64\_nonlinear & 0.530 (ep. 6) & 0.514 (ep. 21) & 0.530 / 0.487 (ep. 6) & 0.037 / 0.505 (ep. 50) \\
infused\_token\_vit\_b\_mlp1\_CORRECTED\_NOSMOOTH & 0.661 (ep. 31) & 0.583 (ep. 50) & 0.657 / 0.583 (ep. 50) & 0.657 / 0.583 (ep. 50) \\
\bottomrule
\end{tabular}
\end{adjustbox}
\caption{Test-set IoU logged during Phase-1 training (the author's wandb runs; validation = the 61 test images, every 5 epochs, input adapter skipped). The saved checkpoint is the epoch with the best mean of TC and AR.}
\label{tab:thesis_recheck_training_logs}
\end{table}


\begin{figure}[htbp]
\centering
\includegraphics[width=\linewidth]{\ResultsDir/figures/thesis_recheck_training_logs.pdf}
\caption{Test IoU logged during the Phase-1 training runs (every five epochs); ring: saved checkpoint.}
\label{fig:recheck_training_logs}
\end{figure}

\subsection{Further checks}

The random-prompt rows cannot be reproduced exactly, because every evaluation draws a new mixture of point and box
prompts (TC spread of 0.03 across three draws). The adapter weights of Table~4.7 are those of the run without label
smoothing, not of the smoothed run whose IoUs Figure~4.2 quotes. The parameter counts of Table~4.11 are exact for all
Infused Token configurations.

\subsection{Phase-2 prompter checkpoints}

The prompter checkpoints of Section~4.3.3 were evaluated with the encoder each was trained with and with their own
decision rule (three-class argmax), Table~\ref{tab:thesis_recheck_phase2} and Figure~\ref{fig:recheck_phase2}. The
multi-scale fusion run whose encoder is unchanged reproduces its logged values exactly (mask TC 0.304 / AR 0.401 versus
0.305 / 0.402). The other generators were trained on an encoder file that was overwritten after the runs; with the
current file their TC IoU is far below the logged values, so they cannot be verified. Where they can be evaluated, SAM
prompted by them is not better than their masks, and the same architectures trained under the protocol of
Section~\ref{sec:prompting_protocol} are equal or better.

\begin{table}[htbp]
\centering
\small
\begin{adjustbox}{max width=\linewidth}
\begin{tabular}{llccccc}
\toprule
Prompter checkpoint (run) & Prompts & Logged mask (best) & Re-run mask & Logged SAM (best) & Re-run SAM & Same arch. retrained \\
\midrule
MSF (Generator) & point & -- / -- & 0.251 / 0.310 & -- / -- & 0.158 / 0.223 & 0.333 / 0.403 \\
MSF (generator\_128\_vit\_b\_bbox) & box & 0.305 / 0.402 & 0.304 / 0.401 & 0.308 / 0.389 & 0.294 / 0.383 & 0.333 / 0.403 \\
MSF + token gate (FIRST) & point & 0.321 / 0.056 & 0.256 / 0.056 & 0.304 / 0.291 & 0.295 / 0.074 & 0.324 / 0.399 \\
token gate + CG blocks (CG) & point & 0.349 / 0.407 & 0.157 / 0.374 & 0.329 / 0.346 & 0.274 / 0.291 & 0.324 / 0.355 \\
token gate + CG blocks (CG\_256) & box & 0.315 / 0.375 & 0.144 / 0.360 & 0.318 / 0.331 & 0.175 / 0.373 & 0.324 / 0.355 \\
token gate + CG blocks (CG\_256\_REAL) & box & 0.327 / 0.387 & 0.118 / 0.352 & 0.322 / 0.332 & 0.141 / 0.371 & 0.324 / 0.355 \\
token gate, shared weights (SP\_TOKEN) & box & 0.326 / 0.393 & 0.205 / 0.376 & 0.333 / 0.306 & 0.190 / 0.361 & 0.332 / 0.402 \\
logistic regression, block 1 (Logistic) & box & 0.191 / 0.295 & 0.190 / 0.287 & 0.183 / 0.255 & 0.194 / 0.298 & 0.195 / 0.319 \\
\bottomrule
\end{tabular}
\end{adjustbox}
\caption{The author's Phase-2 prompter checkpoints (TC IoU / AR IoU, 61 test images), each with the frozen encoder it was trained with. Logged: best value in the author's wandb run (selected on the test set; the run logged SAM with its own prompt pipeline). Re-run mask: 3-class argmax, as in the training scripts. Same arch. retrained: the architecture trained under the benchmark protocol (two sigmoid channels, selection on validation, primary checkpoint, mean of seeds). SAM hybrid and object counts: 08\_thesis\_recheck/phase2\_all.csv.}
\label{tab:thesis_recheck_phase2}
\end{table}


\begin{figure}[htbp]
\centering
\includegraphics[width=\linewidth]{\ResultsDir/figures/thesis_recheck_phase2.pdf}
\caption{The Phase-2 prompter checkpoints of the thesis: best logged value (ring), re-run mask (grey), re-run SAM with
the prompt type used in training (squares), and the same architecture trained under the benchmark protocol (red).}
\label{fig:recheck_phase2}
\end{figure}

\begin{figure}[htbp]
\centering
\includegraphics[width=\linewidth]{\ResultsDir/figures/thesis_recheck_maps_prompters_17.pdf}
\caption{Test image 17: the thesis' multi-scale fusion checkpoint (a) and SAM prompted by its boxes (b), the thesis'
logistic regression on ViT block~1 (c), and multi-scale fusion (d) and the mask-prompt generator (e, f) trained under
the benchmark protocol.}
\label{fig:recheck_maps_prompters}
\end{figure}
